
\uline{sa valeur sur l'élément} @@M0@@ \uline{de} @@M1@@, \uline{est injective}; \uline{son image est le sous-groupe de} @@M2@@ \uline{formé des éléments} @@M3@@ \uline{tels que} (\uline{avec des notations analogues à celles de} 1.4 \uline{et} 1.5)~:
@@M4@@

Le théorème résulte immédiatement du théorème 1.6 et du lemme 1.9, compte tenu de ce que @@M5@@ ne dépend pas de la filtration choisie sur @@M6@@.

Le théorème analogue obtenu en remplaçant @@M7@@ par @@M8@@ est également vrai (cf. 1.7)

Quand il n'y aura pas de risque de confusion, nous noterons le foncteur @@M9@@ par @@M10@@, ce qui est légitime, car @@M11@@ identifie @@M12@@ à une sous-catégorie de @@M13@@.

1.11. On notera @@M14@@ la catégorie des @@M15@@-algèbres graduées, et @@M16@@ le foncteur de @@M17@@ dans @@M18@@ qui à une @@M19@@-algèbre graduée @@M20@@ associe la @@M21@@-algèbre @@M22@@ munie de la filtration canoniquement définie par sa graduation~; le foncteur «~gradué associé~»~: @@M23@@, sera noté @@M24@@. Il existe un isomorphisme de foncteurs @@M25@@.

Soit @@M26@@~; on notera @@M27@@ le groupe multiplicatif des éléments de @@M28@@ d'augmentation @@M29@@. En le munissant de la filtration analogue à celle de @@M30@@, on définit ainsi un foncteur~: @@M31@@.

Par ailleurs, on définit pour tout @@M32@@ un homomorphisme de groupes filtrés~: @@M33@@, en associant à la série @@M34@@, 
